% GENERATED FILE. Edit canonical lessons, then run npm run generate:books.
% canonical-source-hash: fnv1a64:a12d801a852d5c39
% canonical-lessons: RU-C448-prachechnaya, RU-C448-khimchistka, RU-C448-shtopat, RU-C448-zashivat, RU-C448-prishivat

\chapter{Laundry and mending}
\label{ch:ru-laundry-and-mending}
\hlchaptermodality{448}

\begin{chapteropening}
\textbf{By the end of this chapter:} \emph{I can get clothes cleaned and mended: the laundry, the dry cleaner's, darning, sewing up and sewing on.}

Complete the last lesson of chapter 448: I can get clothes cleaned and mended: the laundry, the dry cleaner's, darning, sewing up and sewing on.
\end{chapteropening}

\section[\texorpdfstring{\ru{прачечная} (práchechnaya)}{прачечная (práchechnaya)}]{\ru{прачечная} (práchechnaya) — a laundry (shop), a launderette}

\label{lesson:RU-C448-prachechnaya}



\begin{quote}
Before the new one: say the Russian for a locksmith, then the Russian for repairs.
\end{quote}

\subsection*{You'll want to know: \ru{прачечная}}
\textbf{\ru{прачечная}} (\emph{práchechnaya}) — \textquotedblleft{}a laundry (shop), a launderette\textquotedblright{}.

\begin{quote}
\ru{В} \ru{прачечной} \ru{стирают} \ru{бельё}.
\end{quote}

\begin{quote}
\emph{V práchechnoy stiráyut belyó.}
\end{quote}

\begin{quote}
At the laundry they wash linen.
\end{quote}

An adjective used as a noun. Traditionally \emph{\ru{чн}} is pronounced \emph{\ru{шн}} here.

\begin{cousinweb}
\emph{\ru{прачечная}} is built on \emph{\ru{прачка}} 'a washerwoman' with an adjective ending: the washerwomen's place. EtymDB lists \emph{\ru{прачечная}} beside \emph{\ru{прачка}}.
\end{cousinweb}

\begin{grammarlens}[title={What you've built}]
One more word for telling, explaining and describing at B1.
\end{grammarlens}

\subsection*{Your turn}
\begin{itemize}
\raggedright
  \item \emph{Say it:} \emph{práchechnaya}
  \item \emph{Say it:} \emph{práchechnaya}, once more
  \item \emph{Say it:} a sentence of your own with \emph{práchechnaya}
  \item \emph{Recall:} say the Russian for a locksmith, then the Russian for repairs, then say \emph{práchechnaya} again
\end{itemize}

\begin{tcolorbox}[breakable,colback=teal!4,colframe=teal!35!black,title={Before you move on}]
What does \ru{прачечная} mean? (\textquotedblleft{}A laundry (shop), a launderette\textquotedblright{}.) Say it once more.
\end{tcolorbox}
\section[\texorpdfstring{\ru{химчистка} (khimchístka)}{химчистка (khimchístka)}]{\ru{химчистка} (khimchístka) — a dry cleaner's; dry cleaning}

\label{lesson:RU-C448-khimchistka}



\begin{quote}
Before the new one: say the Russian for repairs, then the Russian for a laundry.
\end{quote}

\subsection*{You'll want to know: \ru{химчистка}}
\textbf{\ru{химчистка}} (\emph{khimchístka}) — \textquotedblleft{}a dry cleaner's; dry cleaning\textquotedblright{}.

\begin{quote}
\ru{Мы} \ru{отдали} \ru{пальто} \ru{в} \ru{химчистку}.
\end{quote}

\begin{quote}
\emph{My otdáli paltó v khimchístku.}
\end{quote}

\begin{quote}
We took the coat to the dry cleaner's.
\end{quote}

Feminine; short for \emph{\ru{химическая} \ru{чистка}} — chemical cleaning.

\begin{cousinweb}
\emph{\ru{химчистка}} is a short form of \emph{\ru{химическая} \ru{чистка}} 'chemical cleaning', from \emph{\ru{химия}} 'chemistry' and \emph{\ru{чистка}} 'cleaning', on \emph{\ru{чистый}} 'clean', which is inherited from Proto-Slavic \emph{čistŭ} (EtymDB).
\end{cousinweb}

\begin{grammarlens}[title={What you've built}]
One more word for telling, explaining and describing at B1.
\end{grammarlens}

\subsection*{Your turn}
\begin{itemize}
\raggedright
  \item \emph{Say it:} \emph{khimchístka}
  \item \emph{Say it:} \emph{khimchístka}, once more
  \item \emph{Say it:} a sentence of your own with \emph{khimchístka}
  \item \emph{Recall:} say the Russian for repairs, then the Russian for a laundry, then say \emph{khimchístka} again
\end{itemize}

\begin{tcolorbox}[breakable,colback=teal!4,colframe=teal!35!black,title={Before you move on}]
What does \ru{химчистка} mean? (\textquotedblleft{}A dry cleaner's; dry cleaning\textquotedblright{}.) Say it once more.
\end{tcolorbox}
\section[\texorpdfstring{\ru{штопать} (shtópat)}{штопать (shtópat)}]{\ru{штопать} (shtópat) — to darn}

\label{lesson:RU-C448-shtopat}



\begin{quote}
Before the new one: say the Russian for a laundry, then the Russian for a dry cleaner's.
\end{quote}

\subsection*{You'll want to know: \ru{штопать}}
\textbf{\ru{штопать}} (\emph{shtópat}) — \textquotedblleft{}to darn\textquotedblright{}.

\begin{quote}
\ru{Бабушка} \ru{штопает} \ru{носки}.
\end{quote}

\begin{quote}
\emph{Bábushka shtópayet noskí.}
\end{quote}

\begin{quote}
Grandma darns socks.
\end{quote}

Imperfective; perfective \emph{\ru{заштопать}}.

\begin{cousinweb}
\emph{\ru{штопать}} was borrowed from Dutch \emph{stoppen} 'to darn' (EtymDB).
\end{cousinweb}

\begin{grammarlens}[title={What you've built}]
One more word for telling, explaining and describing at B1.
\end{grammarlens}

\subsection*{Your turn}
\begin{itemize}
\raggedright
  \item \emph{Say it:} \emph{shtópat}
  \item \emph{Say it:} \emph{shtópat}, once more
  \item \emph{Say it:} a sentence of your own with \emph{shtópat}
  \item \emph{Recall:} say the Russian for a laundry, then the Russian for a dry cleaner's, then say \emph{shtópat} again
\end{itemize}

\begin{tcolorbox}[breakable,colback=teal!4,colframe=teal!35!black,title={Before you move on}]
What does \ru{штопать} mean? (\textquotedblleft{}To darn\textquotedblright{}.) Say it once more.
\end{tcolorbox}
\section[\texorpdfstring{\ru{зашивать} (zashivát)}{зашивать (zashivát)}]{\ru{зашивать} (zashivát) — to sew up, to mend (by sewing)}

\label{lesson:RU-C448-zashivat}



\begin{quote}
Before the new one: say the Russian for a dry cleaner's, then the Russian for to darn.
\end{quote}

\subsection*{You'll want to know: \ru{зашивать}}
\textbf{\ru{зашивать}} (\emph{zashivát}) — \textquotedblleft{}to sew up, to mend (by sewing)\textquotedblright{}.

\begin{quote}
\ru{Мама} \ru{зашивает} \ru{дыру} \ru{на} \ru{брюках}.
\end{quote}

\begin{quote}
\emph{Máma zashiváyet dyrú na bryúkakh.}
\end{quote}

\begin{quote}
Mum is sewing up the hole in the trousers.
\end{quote}

Imperfective; perfective \emph{\ru{зашить}} (\emph{\ru{зашью}, \ru{зашьёшь}}).

\begin{cousinweb}
\emph{\ru{зашивать}} is built from \emph{\ru{за}-} and \emph{-\ru{шивать}}, from \emph{\ru{шить}} 'to sew', which is inherited from Proto-Slavic \emph{šiti} 'to sew' (EtymDB).
\end{cousinweb}

\begin{grammarlens}[title={What you've built}]
One more word for telling, explaining and describing at B1.
\end{grammarlens}

\subsection*{Your turn}
\begin{itemize}
\raggedright
  \item \emph{Say it:} \emph{zashivát}
  \item \emph{Say it:} \emph{zashivát}, once more
  \item \emph{Say it:} a sentence of your own with \emph{zashivát}
  \item \emph{Recall:} say the Russian for a dry cleaner's, then the Russian for to darn, then say \emph{zashivát} again
\end{itemize}

\begin{tcolorbox}[breakable,colback=teal!4,colframe=teal!35!black,title={Before you move on}]
What does \ru{зашивать} mean? (\textquotedblleft{}To sew up, to mend (by sewing)\textquotedblright{}.) Say it once more.
\end{tcolorbox}
\section[\texorpdfstring{\ru{пришивать} (prishivát)}{пришивать (prishivát)}]{\ru{пришивать} (prishivát) — to sew on}

\label{lesson:RU-C448-prishivat}



\begin{quote}
Before the new one: say the Russian for to darn, then the Russian for to sew up.
\end{quote}

\subsection*{You'll want to know: \ru{пришивать}}
\textbf{\ru{пришивать}} (\emph{prishivát}) — \textquotedblleft{}to sew on\textquotedblright{}.

\begin{quote}
\ru{Он} \ru{сам} \ru{пришивает} \ru{пуговицы}.
\end{quote}

\begin{quote}
\emph{On sam prishiváyet púgovitsy.}
\end{quote}

\begin{quote}
He sews on buttons himself.
\end{quote}

Imperfective; perfective \emph{\ru{пришить}}. Takes \emph{\ru{к}} + dative: \emph{\ru{пришить} \ru{пуговицу} \ru{к} \ru{рубашке}}.

\begin{cousinweb}
\emph{\ru{пришивать}} is built from \emph{\ru{при}-} 'onto' and \emph{-\ru{шивать}}, from \emph{\ru{шить}} 'to sew', which is inherited from Proto-Slavic \emph{šiti} (EtymDB).
\end{cousinweb}

\begin{grammarlens}[title={What you've built}]
One more word for telling, explaining and describing at B1.
\end{grammarlens}

\subsection*{Your turn}
\begin{itemize}
\raggedright
  \item \emph{Say it:} \emph{prishivát}
  \item \emph{Say it:} \emph{prishivát}, once more
  \item \emph{Say it:} a sentence of your own with \emph{prishivát}
  \item \emph{Recall:} say the Russian for to darn, then the Russian for to sew up, then say \emph{prishivát} again
\end{itemize}

\begin{tcolorbox}[breakable,colback=teal!4,colframe=teal!35!black,title={Before you move on}]
What does \ru{пришивать} mean? (\textquotedblleft{}To sew on\textquotedblright{}.) Say it once more.
\end{tcolorbox}
